\chapter{Using \textsc{Ply} as a Parser Generator  \label{chapter:ply}}

\begin{figure}[h] 
\centering
  \includegraphics[width=10.5cm]{Abbildungen/golem.png}
\caption{A scientist feeding a grammar to a parser generator golem.}
\label{fig:golem.png}
\end{figure}


Many modern programming languages, e.g.~\texttt{C}, \textsl{Java}, or \textsc{Go} can
be parsed using a so called \simtextsc{Lalr}(1)-Parser\footnote{
  LALR(1) is an acronym that stands for \blue{Look-Ahead Left-to-Right} parsing with 1 token of lookahead.
  We will discuss the theory of \simtextsc{Lalr}(1) parsing at the end of Chapter \ref{chapter:bottom-up}.}.
languages later, this chapter introduces the parser generator
\href{https://www.dabeaz.com/ply/}{\simtextsc{Ply}}, which can be used to 
generate a parser for any language that has an \simtextsc{Lalr}(1) grammar.  In Chapter \ref{chapter:ply-lex} we have already
seen how \simtextsc{Ply} can be used to generate a scanner.  This chapter focuses on the parser-generating aspect
of \simtextsc{Ply}.
We will discuss three example that use \simtextsc{Ply} as a parser generator.
\begin{enumerate}
\item First, we implement a symbolic calculator using \simtextsc{Ply}.
\item In our second example, we develop a program for symbolic differentiation.
      This example shows how \simtextsc{Ply} can be used to construct an abstract syntax tree.
\item In our final example we create an interpreter for a language that is a small fragment of the
      programming language \texttt{C}.
\end{enumerate}


\section{A Symbolic Calculator}
Figure \ref{fig:calculator.g} on page \pageref{fig:calculator.g} shows the grammar of a simple 
\blue{symbolic calculator}, i.e.~a calculator that supports the use of variables.  

\begin{figure}[!ht]

  \begin{center}    
  \framebox{
  \framebox{
  \begin{minipage}[t]{9cm}
  \begin{eqnarray*}
  \textsl{stmnt}   & \rightarrow & \;\simtextsc{Identifier} \quoted{:=} \textsl{expr} \quoted{;}\\
                   & \mid        & \;\textsl{expr} \quoted{;}                   \\[0.2cm]
  \textsl{expr}    & \rightarrow & \;\textsl{expr} \quoted{+} \textsl{product}  \\
                   & \mid        & \;\textsl{expr} \quoted{-} \textsl{product}  \\
                   & \mid        & \;\textsl{product}                           \\[0.2cm]
  \textsl{product} & \rightarrow & \;\textsl{product} \quoted{*} \textsl{factor}\\
                   & \mid        & \;\textsl{product} \quoted{/} \textsl{factor}\\
                   & \mid        & \;\textsl{factor}                            \\[0.2cm]
  \textsl{factor}  & \rightarrow &   \quoted{(} \textsl{expr} \quoted{)}        \\
                   & \mid        & \;\simtextsc{Number}                            \\
                   & \mid        & \;\simtextsc{Identifier}                        
  \end{eqnarray*}
  \vspace*{-0.5cm}
  \end{minipage}}}
  \end{center}
  \caption{A grammar for a symbolic calculator.}
  \label{fig:calculator.g}
\end{figure}

In order to generate a symbolic calculator that is based on this grammar we first need to implement a scanner.
Figure \ref{fig:Symbolic-Calculator.ipynb:lex} shows how to specify an appropriate scanner with \simtextsc{Ply}.
As we have discussed scanner generation with \simtextsc{Ply} at length in Chapter \ref{chapter:ply-lex} there is
no need for further discussions here.

\begin{figure}[!ht]
\centering
\begin{minted}[ frame        = lines, 
                framesep     = 0.3cm, 
                bgcolor      = sepia,
                numbers      = left,
                numbersep    = -0.2cm,
                xleftmargin  = 0.8cm,
                xrightmargin = 0.8cm,
              ]{python3}
    import ply.lex as lex
    
    tokens = [ 'NUMBER', 'IDENTIFIER', 'ASSIGN' ]
    
    def t_NUMBER(t):
        r'0|[1-9][0-9]*(\.[0-9]+)?([eE][+-]?([1-9][0-9]*))?'
        t.value = float(t.value)
        return t
    
    def t_IDENTIFIER(t):
        r'[a-zA-Z][a-zA-Z0-9_]*'
        return t
    
    def t_ASSIGN_OP(t):
        r':='
        return t
    
    literals = ['+', '-', '*', '/', '(', ')', ';']
    
    t_ignore = ' \t'

    def t_newline(t):
        r'\n+'
        t.lexer.lineno += t.value.count('\n')
        
    def t_error(t):
        c = t.value[0]
        l = t.lexer.lineno
        print(f"Illegal character '{c}' in line {l}.")
        t.lexer.skip(1)
    
    lexer = lex.lex()
\end{minted}
\vspace*{-0.3cm}
\caption{A scanner for the symbolic calculator.}
\label{fig:Symbolic-Calculator.ipynb:lex}
\end{figure}

\begin{figure}[!ht]
\centering
\begin{minted}[ frame        = lines, 
                framesep     = 0.3cm, 
                bgcolor      = sepia,
                numbers      = left,
                numbersep    = -0.2cm,
                xleftmargin  = 0.8cm,
                xrightmargin = 0.8cm,
              ]{python3}
    import ply.yacc as yacc
    
    start = 'stmnt'
    
    def p_stmnt_assign(p):
        "stmnt : IDENTIFIER ASSIGN expr ';'"
        Names2Values[p[1]] = p[3]
        
    def p_stmnt_expr(p):
        "stmnt : expr ';'"
        print(p[1])
        
    def p_expr_plus(p):
        "expr : expr '+' prod"
        p[0] = p[1] + p[3]
        
    def p_expr_minus(p):
        "expr : expr '-' prod"
        p[0] = p[1] - p[3]
        
    def p_expr_prod(p):
        "expr : prod"
        p[0] = p[1]
        
    def p_prod_mult(p):
        "prod : prod '*' factor"
        p[0] = p[1] * p[3]
        
    def p_prod_div(p):
        "prod : prod '/' factor"
        p[0] = p[1] / p[3]
        
    def p_prod_factor(p):
        "prod : factor"
        p[0] = p[1]
\end{minted}
\vspace*{-0.3cm}
\caption{A parser for the symbolic calculator, part \texttt{I}.}
\label{fig:Symbolic-Calculator.ipynb:yacc}
\end{figure}

Figure \ref{fig:Symbolic-Calculator.ipynb:yacc} on page \pageref{fig:Symbolic-Calculator.ipynb:yacc} shows how the
grammar is implemented in \simtextsc{Ply}.  We discuss it line by line.
\begin{enumerate}
\item Line 1 imports the module \texttt{ply.yacc}.  This module contains the function \texttt{ply.yacc.yacc}
      which is responsible for computing the parse table. The name \blue{yacc} is a homage to the Unix tool
      \href{https://en.wikipedia.org/wiki/Yacc}{\simtextsc{Yacc}}, which is a popular parser generator for the language
      \texttt{C} and, furthermore, is part of the standard utilities of the Unix operating system.
\item Line 3 specifies that the syntactical variable \texttt{stmnt} is the \blue{start symbol} of the grammar.
\item Line 5 -- 7 define the function \texttt{p\_stmnt\_assign} which implements the grammar rule
      \\[0.2cm]
      \hspace*{1.3cm}
      $\textsl{stmnt} \rightarrow \simtextsc{Identifier} \quoted{:=} \textsl{expr}$.
      \\[0.2cm]
      This grammar rule is represented by the \blue{document string} of the function 
      \texttt{p\_stmnt\_assign}.  In general, if
      \\[0.2cm]
      \hspace*{1.3cm}
      $v \rightarrow \alpha$
      \\[0.2cm]
      is a grammar rule, then this grammar rule is represented by a function that has the name
      \texttt{p\_$v$\_$s$}.  Here, the prefix ``\texttt{p\_}'' specifies that the function implements 
      a grammar rule (the \texttt{p} is short for \blue{parser}), $v$ should\footnote{
        This is just a convention. Technically, $v$ can be any string that is a valid \textsl{Python}
        identifier.
      }
      be the name of the variable
      defined by this grammar rule, and $s$ is a string chosen by the user to distinguish between different
      grammar rules for the same variable.  Of course, $s$ has to be chosen in a way such that the string
      \texttt{p\_$v$\_$s$} is a legal \textsl{Python} identifier.

      The function always takes one argument \texttt{p}.  This argument is a sequence of objects that can be
      indexed with array notation. If the grammar rule defining $v$ has the form
      \\[0.2cm]
      \hspace*{1.3cm}
      $v \rightarrow X_1 \cdots X_n$,
      \\[0.2cm]
      then this sequence has a length of $n+1$.  If $X_i$ is a token, then \texttt{p[$i$]} is the property
      with name \texttt{value} that is associated with this token.  Often, this value is just a string, but it
      can also be a number.  If $X_i$ is a variable, then  \texttt{p[$i$]} is the value that is returned
      when $X_i$ is recognized.  The value associated with the variable $v$ is stored in the location
      \texttt{p[0]}.  In the grammar rule shown in line 5--7, we have not assigned any value to \texttt{p[0]} and
      therefore there is no value associated with the syntactical variable \texttt{stmnt} that is defined by
      this grammar rule.

      \underline{\textbf{Note:}} Line 6 shows how a grammar rule is represented for \simtextsc{Ply}.
      A grammar rule of the form 
      \\[0.2cm]
      \hspace*{1.3cm}
      $v \rightarrow X_1 \cdots X_n$
      \\[0.2cm]
      is represented as the string:
      \\[0.2cm]
      \hspace*{1.3cm}
      \texttt{"$v$ : $X_1$ $\cdots$ $X_n$"}
      \\[0.2cm]
      It is very \colorbox{black}{\red{\underline{im}p\underline{ortant}}} that the character ``\texttt{:}'' is
      surrounded by space characters.  Otherwise, the parser generator does not work but rather generates error
      messages that are difficult to understand. 

      The function \texttt{p\_stmnt\_assign} has the task of evaluating the expression that is on the right
      hand side of the assignment operator ``\texttt{:=}''.  The result of this evaluation is then stored in the
      dictionary \texttt{Names2Values}.  The key that is used is the name of the identifier to the left of the
      assignment operator.
\item The function in line 9 -- 11 implements the grammar rule
      \\[0.2cm]
      \hspace*{1.3cm}
      $\textsl{stmnt} \rightarrow \textsl{expr} \quoted{;}$.
      \\[0.2cm]
      The rule is implemented by evaluating the expression and then printing it.
\item The function \texttt{p\_expr\_plus} implements the grammar rule
      \\[0.2cm]
      \hspace*{1.3cm}
      $\textsl{expr} \rightarrow \textsl{expr} \quoted{+} \textsl{prod}$.
      \\[0.2cm]
      It is implemented by evaluating the expression to the left of the operator ``\texttt{+}'', which is
      stored in \texttt{p[1]}, and the product to the right of this operator, which is stored in \texttt{p[3]},
      and then adding the corresponding values.  Finally, the resulting sum is stored in \texttt{p[0]} so that
      it is available later as the value of the expression that has been parsed.

      The remaining functions are similar to the ones that are discussed above.
\end{enumerate}

\begin{figure}[!ht]
\centering
\begin{minted}[ frame        = lines, 
                framesep     = 0.3cm, 
                bgcolor      = sepia,
                firstnumber  = last,
                numbers      = left,
                numbersep    = -0.2cm,
                xleftmargin  = 0.8cm,
                xrightmargin = 0.8cm,
                ]{python3}
    def p_factor_group(p):
        "factor : '(' expr ')'"
        p[0] = p[2]
        
    def p_factor_number(p):
        "factor : NUMBER"
        p[0] = p[1]

    def p_factor_id(p):
        "factor : IDENTIFIER"
        p[0] = Names2Values.get(p[1], float('nan'))
                
    def p_error(p):
        l = p.lexer.lineno
        if p:
            print(f'Syntax error at {p.value} in line {l}.')
        else:
            print('Syntax error at end of input.')
        
    parser = yacc.yacc()

    Names2Values = {}
    
    def main():
        while True:
            s = input('calc > ')
            if s == '':
                break
            yacc.parse(s)
\end{minted}
\vspace*{-0.3cm}
\caption{A parser for the symbolic calculator, part \texttt{II}.}
\label{fig:Symbolic-Calculator.ipynb:yacc2}
\end{figure}
\FloatBarrier

\noindent
Figure \ref{fig:Symbolic-Calculator.ipynb:yacc2} on page \pageref{fig:Symbolic-Calculator.ipynb:yacc2}
is discussed next.
\begin{enumerate}
\item Line 48 -- 43 shows the function \texttt{p\_error} which is used to print error messages in the case that
      the input can not be parsed because of a syntax error.  The argument \texttt{p} is the token $t$ that
      caused the entry $\texttt{action}(s, t)$ in the action table to be undefined.
      If the syntax error happens at the end of the input, \texttt{p} has the value \texttt{None}.
\item Line 55 generates the parser.
\item Line 57 initializes the dictionary \texttt{Names2Values}.  For every identifier $x$ defined interactively,
      \texttt{Names2Values[$x$]} is the value associated with $x$.
\item The function \texttt{main} is used as a driver for the parser.  It reads a string \texttt{s}
      from the command line and tries to parse \texttt{s} using the function \texttt{yacc.parse}.
      The function \texttt{yacc.parse} is generated behind the scenes when the function \texttt{yacc.yacc} is
      invoked in line 55. 
\end{enumerate}

The Jupyter notebook \texttt{01-Calculator}, which is available at
\\[0.2cm]
\hspace*{-0.3cm}
\href{https://github.com/karlstroetmann/Formal-Languages/blob/master/Python/Chapter-07/01-Calculator.ipynb}{https://github.com/karlstroetmann/Formal-Languages/blob/master/Python/Chapter-07/01-Calculator.ipynb}
\\[0.2cm]
implements the symbolic calculator discussed in this section.

\section{Symbolic Differentiation}
Our next example shows how we can compute an abstract syntax tree with the help of \simtextsc{Ply}.
The Jupyter notebook \texttt{02-Differentiator}, which is available at
\\[0.2cm]
\hspace*{-0.3cm}
\href{https://github.com/karlstroetmann/Formal-Languages/blob/master/Python/Chapter-07/02-Differentiator.ipynb}{https://github.com/karlstroetmann/Formal-Languages/blob/master/Python/Chapter-07/02-Differentiator.ipynb}
\\[0.2cm]
implements a parser that generates an abstract syntax tree, which can then be used for symbolic
differentiation.

\section{Implementing a Simple Interpreter  \label{chapter:interpreter}}
Figure \ref{fig:Pure.g} show the grammar of a simple programming language.  In this section, we will develop an
interpreter for this language.

\begin{figure}[!ht]
\centering
\begin{Verbatim}[ frame         = lines, 
                  framesep      = 0.3cm, 
                  firstnumber   = 1,
                  labelposition = bottomline,
                  numbers       = left,
                  numbersep     = -0.2cm,
                  xleftmargin   = 0.8cm,
                  xrightmargin  = 0.8cm,
                ]
    program
        : /* lambda */
        | stmnt program
        
    stmnt 
        : IF '(' bool_expr ')' stmnt                 
        | WHILE '(' bool_expr ')' stmnt
        | '{' program '}' 
        | ID ':=' expr ';'  
        | expr ';'       
    
    bool_expr 
        : expr '==' expr     
        | expr '!=' expr     
        | expr '<=' expr     
        | expr '>=' expr     
        | expr '<'  expr      
        | expr '>'  expr     
     
    expr: expr '+' product                 
        | expr '-' product
        | product
                  
    product
        : product '*' factor               
        | product '/' factor 
        | factor
    
    factor
        : '(' expr ')' 
        | NUMBER
        | ID
        | ID '(' expr_list ')'
    
    expr_list
        : /* lambda */
        | expr ',' ne_expr_list
    
    ne_expr_list
        : expr
        | expr ',' ne_expr_list
\end{Verbatim}
\vspace*{-0.3cm}
\caption{A grammar for a simple programming language.}
\label{fig:Pure.g}
\end{figure}

The statements of the language that we want to implement include \textsl{assignments}, \textsl{function calls},
\textsl{\texttt{if} statements}, and \textsl{\texttt{while} loops}.  The language supports boolean expression
that can be build with the help of the comparison operators
\\[0.2cm]
\hspace*{1.3cm}
\texttt{==}, \texttt{!=}, \texttt{<=}, \texttt{>=}, \texttt{<}, and \texttt{>}.
\\[0.2cm]
Arithmetic expression can use the arithmetic operators 
\\[0.2cm]
\hspace*{1.3cm}
\texttt{+}, \texttt{+}, \texttt{-}, \texttt{*}, and \texttt{/}.
\\[0.2cm]
In the base case, arithmetic expressions are build from numbers, variables, and function calls.
Although the language supports the calling of predefined functions like \texttt{read} and \texttt{print},
it does not support user defined functions.

Figure \ref{fig:sum.sl} shows
an example program that conforms to this grammar. This program first reads a number
which is stored in the variable \texttt{n}. Subsequently, the sum
\\[0.2cm]
\hspace*{1.3cm}
$\ds s := \sum\limits_{\mbox{$\normalsize i=1$}}^{\mbox{$\normalsize n^2$}} \mbox{\Large $i$}$
\\[0.2cm]
is computed using a \texttt{while} loop.  Finally, this sum is printed.

\begin{figure}[!ht]
\centering
\begin{Verbatim}[ frame         = lines, 
                  framesep      = 0.3cm, 
                  labelposition = bottomline,
                  numbers       = left,
                  numbersep     = -0.2cm,
                  xleftmargin   = 0.8cm,
                  xrightmargin  = 0.8cm,
                ]
    n := read();
    s := 0;
    i := 0;
    while (i < n * n) {
        i := i + 1;
        s := s + i;
    }
    print(s);
\end{Verbatim}
\vspace*{-0.3cm}
\caption{A program to compute the sum $\sum\limits_{i=0}^{n^2} i$.}
\label{fig:sum.sl}
\end{figure}
\FloatBarrier

Similar to our program for symbolic differentiation, we will represent the individual commands as
nested tuples. The program shown in Figure \ref{fig:sum.sl} is represented by the nested tuple displayed in
Figure \ref{fig:sum.ast}. This nested tuple is nothing other than the abstract syntax tree
corresponding to the original program. Note that lists of commands are represented as nested tuples
that start with a ``\texttt{.}''.

\begin{figure}[!ht]
\centering
\begin{Verbatim}[ frame         = lines, 
                  framesep      = 0.3cm, 
                  labelposition = bottomline,
                  numbers       = left,
                  numbersep     = -0.2cm,
                  xleftmargin   = 0.8cm,
                  xrightmargin  = 0.8cm,
                ]
    ('.',            
        ('read', 'n'),
        (':=', 's', 0),
        (':=', 'i', 0),
        ('while', ('<',  'i', ('*', 'n', 'n')),
            ('.',
                (':=', 'i', ('+', 'i', 1)),
                (':=', 's', ('+', 's', 'i'))
            )    
        ),
        ('print', 's')
    )
\end{Verbatim}
\vspace*{-0.3cm}
\caption{A nested tuple that represents the program from Figure \ref{fig:sum.sl}.}
\label{fig:sum.ast}
\end{figure}

\begin{figure}[!ht]
\centering
\begin{minted}[ frame         = lines, 
                framesep      = 0.3cm,
                bgcolor       = sepia,
                firstnumber   = 1,
                numbers       = left,
                numbersep     = -0.3cm,
                xleftmargin   = 0.3cm,
                xrightmargin  = 0.0cm,
              ]{python3}
    tokens = [ 'NUMBER',    # r'0|[1-9][0-9]*'
               'ID',        # r'[a-zA-Z][a-zA-Z0-9_]*'
               'ASSIGN',    # r':='
               'EQ',        # r'=='
               'NE',        # r'!='
               'LE',        # r'<='
               'GE',        # r'>='
               'IF',        # r'if'
               'WHILE'      # r'while'
              ]
    
    def t_COMMENT(t):
        r'(/\*(.|\n)*?\*/)|(//.*)'
        t.lexer.lineno += t.value.count('\n')
        pass
    
    def t_NUMBER(t):
        r'0|[1-9][0-9]*'
        t.value = int(t.value)
        return t
    
    t_ASSIGN = r':='
    t_EQ     = r'=='
    t_NE     = r'!='
    t_LE     = r'<='
    t_GE     = r'>='
    
    Keywords = { 'if': 'IF', 'while': 'WHILE' }
    
    def t_IDENTIFIER(t):
        r'[a-zA-Z][a-zA-Z0-9_]*'
        t.type = Keywords.get(t.value, 'ID')
        return t
    
    literals = ['+', '-', '*', '/', '%', '(', ')', '{', '}', ';', '<', '>', ',']
    
    t_ignore  = ' \t\r'
\end{minted}
\vspace*{-0.3cm} %\$
\caption{Implementation of the scanner.}
\label{fig:Interpreter-Scanner}
\end{figure}

Figure \ref{fig:Interpreter-Scanner} on \pageref{fig:Interpreter-Scanner} show the implementation of the
scanner. The scanner primarily differentiates between variables (token \simtextsc{Identifier}) and numbers (token
\simtextsc{Number}).  Variables start with an uppercase or lowercase letter, followed by additional digits and 
the underscore. Furthermore, there are operator symbols for the comparison operators
``\texttt{==}'', ``\texttt{!=}'', ``\texttt{<=}'', ``\texttt{>=}'', ``\texttt{<}'', and ``\texttt{>}'',
for the arithmetic operators ``\texttt{+}'',  ``\texttt{-}'',  ``\texttt{*}'',  ``\texttt{/}'', the parentheses
``\texttt{(}'' and ``\texttt{)}'',  and for the assignment operator ``\texttt{:=}''.  Note that only those operators names 
that consist of more than two characters have to be declared as tokens, since the other symbols can be declared
as \texttt{literals}.  Finally, the language is equipped with two keywords ``\texttt{if}'' and ``\texttt{while}''.
The scanner is pretty standard, but there are two things worth mentioning.
\begin{enumerate}
\item The scanner removes white space and comments.  The language supports two kinds of comments:
      \begin{enumerate}
      \item The language supports multiline comments, i.e.~text surrounded by ``\texttt{/*}'' ``\texttt{*/}''.
            These comments have the form
            \\[0.2cm]
            \hspace*{1.3cm}
            \texttt{/* ... */}
            \\[0.2cm]
            where ``\texttt{...}'' denotes arbitrary text not containing the substring ``\texttt{*/}''.

            We have to use the so-called \emph{non-greedy} version of the
            Kleene operator ``\texttt{*}'' in the specification of multiline comments. The non-greedy version of the
            operator ``\texttt{*}'' is written as ``\texttt{*?}'' and matches as little as possible. Therefore, the
            regular expression 
            \\[0.2cm]
            \hspace*{1.3cm}
            \verb$/\*(.|\n)*?\*/$ %$
            \\[0.2cm]
            represents a string that starts with the character sequence ``\texttt{/*}'', ends with the
            character sequence ``\texttt{*/}'', and is as short as possible. This ensures that in a line like
            \\[0.2cm]
            \hspace*{1.3cm}
            \texttt{/* Hugo */ i := i + 1; /* Anton */}
            \\[0.2cm]
            two separate comments are recognized.
      \item Furthermore, the language supports one line comments that start with the character sequence ``\texttt{//}''
            and extend to the end of the line.
      \end{enumerate}
      Note that the symbol \simtextsc{Comment} is not declared as a token.  The reason is that it won't appear in the grammar
      rules later, as comments are discarded by the scanner.
\item The treatment of keywords requires special care.  Since both ``\texttt{if}'' and ``\texttt{while}'' have
      the same syntactical form as variables, the function \texttt{t\_IDENTIFIER} recognizes them.  This is done with
      the help of the dictionary \texttt{Keywords}.  If \texttt{t\_IDENTIFIER} finds a string of characters, it first
      checks, whether this string has an associated token value in the dictionary \texttt{Keywords}.  If there
      is a value associated with the string, the string is a keyword and the corresponding token is returned.
      On the other hand, if the string is not stored in the dictionary \texttt{Keywords}, the function
      \texttt{get} returns the token \simtextsc{Identifier} as a default. 
\end{enumerate}


\noindent
The figures \ref{fig:Interpreter-Parser-1}, \ref{fig:Interpreter-Parser-2}, \ref{fig:Interpreter-Parser-3a},
\ref{fig:Interpreter-Parser-3b}, and \ref{fig:Interpreter-Parser-4} show the implementation of the parser using
\simtextsc{Ply}.  Figure \ref{fig:Interpreter-Parser-1} on page \pageref{fig:Interpreter-Parser-1} implements
the following grammar rules: 

\begin{Verbatim}[ frame         = lines, 
                  framesep      = 0.3cm, 
                  firstnumber   = 1,
                  labelposition = bottomline,
                  numbers       = left,
                  numbersep     = -0.2cm,
                  xleftmargin   = 0.8cm,
                  xrightmargin  = 0.8cm,
                ]
    program
        : /* epsilon */
        | stmnt program
        
    stmnt 
        : IF '(' bool_expr ')' stmnt                 
        | WHILE '(' bool_expr ')' stmnt
        | '{' program '}' 
        | IDENTIFIER ':=' expr ';'  
        | expr ';'       
\end{Verbatim}



\begin{figure}[!ht]
\centering
\begin{minted}[ frame         = lines, 
                framesep      = 0.3cm, 
                labelposition = bottomline,
                numbers       = left,
                numbersep     = -0.3cm,
                xleftmargin   = 0.3cm,
                xrightmargin  = 0.0cm,
              ]{python3}
    def p_program_lambda(p):
        "program : "
        p[0] = ('.',)
    
    def p_program_more(p):
        "program : stmnt program"
        p[0] = ('.', p[1]) + p[2][1:]
    
    def p_stmnt_if(p):
        "stmnt : IF '(' bool_expr ')' stmnt"
        p[0] = ('if', p[3], p[5])   
    
    def p_stmnt_while(p):
        "stmnt : WHILE '(' bool_expr ')' stmnt"
        p[0] = ('while', p[3], p[5])
        
    def p_stmnt_block(p):
        "stmnt : '{' program '}'"
        p[0] = p[2]
        
    def p_stmnt_assign(p):
        "stmnt : IDENTIFIER ASSIGN expr ';'"
        p[0] = (':=', p[1], p[3])
       
    def p_stmnt_expr(p):
        "stmnt : expr ';'"
        p[0] = ('expr', p[1])
\end{minted}
\vspace*{-0.3cm} %\$
\caption{Specification of statements.}
\label{fig:Interpreter-Parser-1}
\end{figure}


The start symbol of the grammar is the variable \texttt{program}.
When parsing this variable, the parser returns a nested tuple of statements.
The first element of this tuple is a dot ``\texttt{.}''.

The syntactic variable \texttt{stmnt} describes the various commands that 
are supported in our simple language.
\begin{enumerate}[(a)]
\item A conditional test has the syntax:
      \\[0.2cm]
      \hspace*{1.3cm}
      $\texttt{if}\; \texttt{(}\; b\; \texttt{)}\; \textsl{stmnt}\; $
      \\[0.2cm]
      Here, $b$ is an expression that evaluates to \texttt{True} or \texttt{False}
      and \texttt{stmnt} is a statement that is executed if
      $b$ evaluates to \texttt{True}.
      This command is represented by the nested tuple
      \\[0.2cm]
      \hspace*{1.3cm}
      $(\texttt{'if'}, b, \textsl{stmnt})$.
      \\[0.2cm]
      Here, $b$ is a nested tuple representing the Boolean expression.
\item A loop has the syntax:
      \\[0.2cm]
      \hspace*{1.3cm}
      $\texttt{while}\; \texttt{(}\; b\; \texttt{)}\;  \textsl{stmnt}$.
      \\[0.2cm]
      Here, $b$ is an expression that evaluates to \texttt{True} or \texttt{False}
      and \texttt{stmnt} is a statement that is executed as long as
      $b$ evaluates to \texttt{True}.
      This statement is represented by the nested tuple
      \\[0.2cm]
      \hspace*{1.3cm}
      $(\texttt{'while'}, b, \textsl{stmnt})$.
\item A block statement is a list of statements that is surrounded by curly braces.
\item The simplest commands are the assignments. These have the form:
      \\[0.2cm]
      \hspace*{1.3cm}
      $v \;\texttt{:=}\; e \texttt{;}$
      \\[0.2cm]
      Here, $v$ is a variable and $e$ is an arithmetic expression.
      An assignment is represented by the nested tuple
      \\[0.2cm]
      \hspace*{1.3cm}
      $(\texttt{':='}, v, e)$.
\item Furthermore, every expression is a statement if it is followed by a semicolon.
      The reason to allow this is that some expressions have side effects.  Syntactically,
      the string ``\texttt{print(1+2)}'' is an expression as it is a function call.
      By writing it as
      \\[0.2cm]
      \hspace*{1.3cm}
      \texttt{print(1+2);}
      \\[0.2cm]
      we can execute it as a statement.
\end{enumerate}


\begin{figure}[!ht]
\centering
\begin{minted}[ frame         = lines, 
                framesep      = 0.3cm, 
                labelposition = bottomline,
                numbers       = left,
                numbersep     = -0.3cm,
                xleftmargin   = 0.3cm,
                xrightmargin  = 0.0cm,
              ]{python3}
    def p_bool_expr_eq(p):
        "bool_expr : expr EQ expr"
        p[0] = ('==', p[1], p[3])
    
    def p_bool_expr_ne(p):
        "bool_expr : expr NE expr"
        p[0] = ('!=', p[1], p[3])
    
    def p_bool_expr_le(p):
        "bool_expr : expr LE expr"
        p[0] = ('<=', p[1], p[3])
        
    def p_bool_expr_ge(p):
        "bool_expr : expr GE expr"
        p[0] = ('>=', p[1], p[3])
        
    def p_bool_expr_lt(p):
        "bool_expr : expr '<' expr"
        p[0] = ('<', p[1], p[3])
    
    def p_bool_expr_gt(p):
        "bool_expr : expr '>' expr"
        p[0] = ('>', p[1], p[3])
    
\end{minted}
\vspace*{-0.3cm} %\$
\caption{Specification of boolean expressions.}
\label{fig:Interpreter-Parser-2}
\end{figure}

Figure \ref{fig:Interpreter-Parser-2} shows the parsing of boolean expressions.  It implements the following
grammar rule.
\pagebreak

\begin{Verbatim}[ frame         = lines, 
                  framesep      = 0.3cm, 
                  firstnumber   = 1,
                  labelposition = bottomline,
                  numbers       = left,
                  numbersep     = -0.2cm,
                  xleftmargin   = 0.8cm,
                  xrightmargin  = 0.8cm,
                ]
    bool_expr 
        : expr '==' expr     
        | expr '!=' expr     
        | expr '<=' expr     
        | expr '>=' expr     
        | expr '<'  expr      
        | expr '>'  expr     
\end{Verbatim}

The syntactic variable \texttt{boolExpr} describes an expression that takes a Boolean value.
\begin{enumerate}
\item An expression of the form
      \\[0.2cm]
      \hspace*{1.3cm}
      $l \;\texttt{==}\; r $
      \\[0.2cm]
      tests whether the evaluations of $l$ and $r$ yield the same values.
      This expression is represented by the nested tuple
      \\[0.2cm]
      \hspace*{1.3cm}
      $(\texttt{'=='}, l, r)$.
\item The other cases are similar.
\end{enumerate}


The figure \ref{fig:Interpreter-Parser-3a} and \ref{fig:Interpreter-Parser-3b} on page
\pageref{fig:Interpreter-Parser-3a} implements the following grammar rules:
\begin{Verbatim}[ frame         = lines, 
                  framesep      = 0.3cm, 
                  firstnumber   = 1,
                  labelposition = bottomline,
                  numbers       = left,
                  numbersep     = -0.2cm,
                  xleftmargin   = 0.8cm,
                  xrightmargin  = 0.8cm,
                ]
    expr: expr '+' product                 
        | expr '-' product
        | product
                  
    product
        : product '*' factor               
        | product '/' factor
        | product '%' factor 
        | factor
    
    factor
        : '(' expr ')' 
        | NUMBER
        | IDENTIFIER
        | IDENTIFIER '(' expr_list ')'
\end{Verbatim}


\begin{figure}[!ht]
\centering
\begin{minted}[ frame         = lines, 
                framesep      = 0.3cm, 
                labelposition = bottomline,
                numbers       = left,
                numbersep     = -0.3cm,
                xleftmargin   = 0.3cm,
                xrightmargin  = 0.0cm,
              ]{python3}
    def p_expr_plus(p):
        "expr : expr '+' product"
        p[0] = ('+', p[1], p[3])        
    def p_expr_minus(p):
        "expr : expr '-' product"
        p[0] = ('-', p[1], p[3])
    
    def p_expr_product(p):
        "expr : product"
        p[0] = p[1]
        
    def p_product_times(p):
        "product : product '*' expr"
        p[0] = ('*', p[1], p[3])
    def p_product_divide(p):
        "product : product '/' expr"
        p[0] = ('/', p[1], p[3])
    
    def p_product_modulo(p):
        "product : product '%' factor"
        p[0] = ('%', p[1], p[3])
    
    def p_product_factor(p):
        "product : factor"
        p[0] = p[1]
\end{minted}
\vspace*{-0.3cm} %\$
\caption{Specification of arithmetic expressions, part \texttt{I}.}
\label{fig:Interpreter-Parser-3a}
\end{figure}


\begin{figure}[!ht]
\centering
\begin{minted}[ frame         = lines, 
                framesep      = 0.3cm, 
                firstnumber   = last,
                labelposition = bottomline,
                numbers       = left,
                numbersep     = -0.3cm,
                xleftmargin   = 0.3cm,
                xrightmargin  = 0.0cm,
              ]{python3}
    def p_product_factor(p):
        "product : factor"
        p[0] = p[1]
    
    def p_factor_paren(p):
        "factor : '(' expr ')'"
        p[0] = p[2]
    
    def p_factor_number(p):
        "factor : NUMBER"
        p[0] = p[1]
    
    def p_factor_id(p):
        "factor : IDENTIFIER"
        p[0] = p[1]
    
    def p_factor_fct_call(p):
        "factor : IDENTIFIER '(' expr_list ')'"
        p[0] = ('call', p[1]) + p[3][1:]
\end{minted}
\vspace*{-0.3cm} %\$
\caption{Specification of arithmetic expressions, part \texttt{II}.}
\label{fig:Interpreter-Parser-3b}
\end{figure}

\noindent
To understand the code shown in Figure \ref{fig:Interpreter-Parser-3a} and \ref{fig:Interpreter-Parser-3b} we
have to understand the following: 
\begin{enumerate}
\item An arithmetic expression of the form $a + b$ is represented as a triple of the form
      \\[0.2cm]
      \hspace*{1.3cm}
      \texttt{('+', x, y)}
      \\[0.2cm]
      where \texttt{x} and \texttt{y} are the tuples corresponding to the abstract syntax trees of $a$ and $b$,
      respectively. 
\item Arithmetic expressions using the operators ``\texttt{-}'', ``\texttt{*}'', ``\texttt{/}'', ``\texttt{\%}''
      are represented analogously.
\item A parenthesized expression, i.e.~an expression of the form
      \\[0.2cm]
      \hspace*{1.3cm}
      $(e)$
      \\[0.2cm]
      is represented by the abstract tree that represents $e$.  Therefore, the parentheses are just used for
      grouping and do not appear in the abstract syntax tree.
\item A number is represented by itself.
\item A variable is represented by its name, i.e.~the variable \texttt{x} is represented by the string
      ``\texttt{x}''.
\item A function call of the form
      \\[0.2cm]
      \hspace*{1.3cm}
      $f(a_1, \cdots a_n)$
      \\[0.2cm]
      is represented by the tuple
      \\[0.2cm]
      \hspace*{1.3cm}
      \texttt{('call', $f$, $a_1$, $\cdots$, $a_n$)}.
\end{enumerate}
\pagebreak

\noindent
Figure \ref{fig:Interpreter-Parser-4} on page \pageref{fig:Interpreter-Parser-4} implements the following
grammar rules:
\begin{Verbatim}[ frame         = lines, 
                  framesep      = 0.3cm, 
                  firstnumber   = 1,
                  labelposition = bottomline,
                  numbers       = left,
                  numbersep     = -0.2cm,
                  xleftmargin   = 0.8cm,
                  xrightmargin  = 0.8cm,
                ]
    expr_list
        :
        | expr ',' ne_expr_list
    
    ne_expr_list
        : expr
        | expr ',' ne_expr_list
\end{Verbatim}


\begin{figure}[!ht]
\centering
\begin{minted}[ frame         = lines, 
                framesep      = 0.3cm, 
                labelposition = bottomline,
                numbers       = left,
                numbersep     = -0.3cm,
                xleftmargin   = 0.3cm,
                xrightmargin  = 0.0cm,
              ]{python3}
    def p_expr_list_empty(p):
        "expr_list : "
        p[0] = ('.',)
        
    def p_expr_list_one(p):
        "expr_list : expr"
        p[0] = ('.', p[1])     
    
    def p_expr_list_more(p):
        "expr_list : expr ',' ne_expr_list"
        p[0] = ('.', p[1]) + p[3][1:]     
    
    def p_ne_expr_list_one(p):
        "ne_expr_list : expr"
        p[0] = ('.', p[1]) 
        
    def p_ne_expr_list_more(p):
        "ne_expr_list : expr ',' ne_expr_list"
        p[0] = ('.', p[1]) + p[3][1:] 
\end{minted}
\vspace*{-0.3cm} %\$
\caption{Specification of lists of expressions.}
\label{fig:Interpreter-Parser-4}
\end{figure}




\begin{figure}[!ht]
\centering
\begin{minted}[ frame         = lines, 
                framesep      = 0.3cm, 
                labelposition = bottomline,
                numbers       = left,
                numbersep     = -0.3cm,
                xleftmargin   = 0.3cm,
                xrightmargin  = 0.0cm,
              ]{python3}
    def main(file):
        with open(file, 'r') as handle:
            program = handle.read() 
        stmnt = yacc.parse(program)
        Values = {}
        execute(stmnt, Values)
\end{minted}
\vspace*{-0.3cm}
\caption{Specification of the function \texttt{main}.}
\label{fig:Interpreter.ipynb:main}
\end{figure}

\begin{figure}[!ht]
\centering
\begin{minted}[ frame         = lines, 
                framesep      = 0.3cm, 
                labelposition = bottomline,
                numbers       = left,
                numbersep     = -0.3cm,
                xleftmargin   = 0.3cm,
                xrightmargin  = 0.0cm,
              ]{python3}
    def execute(stmnt: NestedTuple, Values: dict[str, Number]) -> None:
        match stmnt:
            case ('.', *SL):
                execute_tuple(tuple(SL), Values)
            case (':=', var, value):
                Values[var] = evaluate(value, Values)
            case ('expr', expr):
                evaluate(expr, Values)
            case ('if', test, stmnt):
                if evaluate_bool(test, Values):
                    execute(stmnt, Values)
            case ('while', test, stmnt):
                while evaluate_bool(test, Values):
                    execute(stmnt, Values)
            case _:
                assert False, f'{stmnt} unexpected'

    def execute_tuple(Statement_List, Values={}):
        for stmnt in Statement_List:
            execute(stmnt, Values)
\end{minted}
\vspace*{-0.3cm}
\caption{The function \texttt{execute}.}
\label{fig:Interpreter.ipynb:execute}
\end{figure}

Figure \ref{fig:Interpreter.ipynb:main} shows the function \texttt{main}, which receives the name of a
file containing a program in our simple programming language as input. This program is parsed and
thus transformed into a tuple of statements. The function \texttt{execute} executes this tuple.  In order to do
so, it needs to know the values of all variables.  Theses are stored in the dictionary \texttt{Values}, which
initially is empty.  Every time a variable is assigned, the variable and the corresponding value will be
stored in this dictionary.

Figure \ref{fig:Interpreter.ipynb:execute} shows the implementation of the function \texttt{execute}. This
implementation consists mainly of a large case distinction based on the type of command to be executed.
\begin{enumerate}
\item First, we check if \texttt{stmnt} is a list of statements. A statement is a list of statements
      if the first element of the tuple representing this statement is the string \texttt{'.'}.

      In this case, the statements following the string \texttt{'.'} are executed via the auxiliary function
      \texttt{execute\_tuple}. 
\item If the statement is an assignment of the form
      \\[0.2cm]
      \hspace*{1.3cm}
      \texttt{(':=', var, value)}
      \\[0.2cm]   
      then the value of the expression \texttt{value} is computed using the function \texttt{evaluate}.
      This value is then stored in the dictionary \texttt{Values} under the key \texttt{var}.
\item If the statement is an expression statement of the form
      \\[0.2cm]
      \hspace*{1.3cm}
      \texttt{(expr', expr)}
      \\[0.2cm]
      then the corresponding \texttt{expr} is evaluated using the function \texttt{evaluate}.
\item If the command is of the form
      \\[0.2cm]
      \hspace*{1.3cm}
      \texttt{('if', test, stmnt)}
      \\[0.2cm]
      then \texttt{test} is a nested tuple representing a Boolean expression, while \texttt{stmnt} is a statement.
      In this case, the expression \texttt{test} is first evaluated using the function \texttt{evaluate}.
      If this evaluation yields the value \texttt{True}, then \texttt{stmnt} is executed.
\item If the command is of the form
      \\[0.2cm]
      \hspace*{1.3cm}
      \texttt{('while', test, stmnt)}
      \\[0.2cm]
      then \texttt{test} is a Boolean expression and \texttt{stmnt} is a statement.
      In this case, the expression \texttt{test} is first evaluated using the function \texttt{evaluate}.
      If this evaluation yields the value \texttt{True}, then \texttt{stmnt} is executed.
      Then the expression \texttt{test} is evaluated again. If the result is
      \texttt{False}, then the execution of the while loop finishes.  Otherwise,
      \texttt{stmnt} is executed as long as the evaluation of \texttt{test}
      yields \texttt{True}.
\end{enumerate}


\begin{figure}[!ht]
\centering
\begin{minted}[ frame         = lines, 
                framesep      = 0.3cm, 
                labelposition = bottomline,
                numbers       = left,
                numbersep     = -0.3cm,
                xleftmargin   = 0.3cm,
                xrightmargin  = 0.0cm,
              ]{python3}
    def evaluate_bool(expr, Values):
        match expr:
            case ('==', lhs, rhs):
                return evaluate(lhs, Values) == evaluate(rhs, Values)
            case ('!=', lhs, rhs):
                return evaluate(lhs, Values) != evaluate(rhs, Values)
            case ('<=', lhs, rhs):
                return evaluate(lhs, Values) <= evaluate(rhs, Values)
            case ('>=', lhs, rhs):
                return evaluate(lhs, Values) >= evaluate(rhs, Values)
            case ('<', lhs, rhs):
                return evaluate(lhs, Values) <  evaluate(rhs, Values)
            case ('>', lhs, rhs):
                return evaluate(lhs, Values) >  evaluate(rhs, Values)
            case _:
                assert False, f'{expr} unexpected'
\end{minted}
\vspace*{-0.3cm}
\caption{The evaluation of boolean expressions.}
\label{fig:Interpreter.ipynb:evaluate_bool}
\end{figure}

Figure \ref{fig:Interpreter.ipynb:evaluate_bool} shows the implementation of the function \texttt{evaluate\_bool}.
This function receives a boolean expression and a dictionary storing the current values of the variables as input.
\begin{enumerate}[(a)]
\item If the expression to be evaluated is a Boolean expression of the form
      \\[0.2cm]
      \hspace*{1.3cm}
      \texttt{('==', lhs, rhs)}
      \\[0.2cm]
      we recursively evaluate the expressions \texttt{lhs} and \texttt{rhs} and return \texttt{True} if and only if
      both expressions yield the same value.
\item If the expression to be evaluated is a Boolean expression of the form
      \\[0.2cm]
      \hspace*{1.3cm}
      \texttt{('<', lhs, rhs)}
      \\[0.2cm]
      we recursively evaluate the expressions \texttt{lhs} and \texttt{rhs} and return \texttt{True} if and only if
      the value obtained from the evaluation of \texttt{lhs} is less than the
      value obtained from the evaluation of \texttt{rhs}.

      The remaining cases are similar.
\end{enumerate}

\begin{figure}[!ht]
\centering
\begin{minted}[ frame         = lines, 
                framesep      = 0.3cm, 
                labelposition = bottomline,
                numbers       = left,
                numbersep     = -0.3cm,
                xleftmargin   = 0.3cm,
                xrightmargin  = 0.0cm,
              ]{python3}
    def evaluate(expr, Values):
        match expr:
            case int():
                return expr
            case str():
                return Values[expr] 
            case ('call', 'read'):
                return int(input('Please enter a natural number: '))
            case ('call', 'print', expr):
                print(evaluate(expr, Values))        
                return 0;
            case ('+', lhs, rhs):
                return evaluate(lhs, Values) + evaluate(rhs, Values)
            case ('-', lhs, rhs):
                return evaluate(lhs, Values) - evaluate(rhs, Values)
            case ('*', lhs, rhs):
                return evaluate(lhs, Values) * evaluate(rhs, Values)
            case ('/', lhs, rhs):
                return evaluate(lhs, Values) / evaluate(rhs, Values)
            case ('%', lhs, rhs):
                return evaluate(lhs, Values) % evaluate(rhs, Values)
            case _:
                assert False, f'{expr} unexpected'
\end{minted}
\vspace*{-0.3cm}
\caption{The evaluation of arithmetic expressions.}
\label{fig:Interpreter.ipynb:evaluate}
\end{figure}

Figure \ref{fig:Interpreter.ipynb:evaluate} shows the implementation of the function \texttt{evaluate}.
This function receives an arithmetic expression and a dictionary storing the
values of the variables as input.
\begin{enumerate}[(a)]
\item If the expression to be evaluated is a number, we return this number as the result.
\item If the expression to be evaluated is a variable, we look up the value of this
      variable in the dictionary \texttt{Values} and return this value as the result.
\item If the expression to be evaluated is a call to the function \texttt{read},
      we prompt the user to enter a natural number. We then convert the string entered by the user into an integer.
\item If the expression to be evaluated is a call to the function \texttt{print}, we evaluate the argument and
      print it.  Furthermore, we return 0 since the evaluation of an expression always has to return a value.
\item If the expression to be evaluated is a sum of the form
      \\[0.2cm]
      \hspace*{1.3cm}
      \texttt{('+', lhs, rhs)}
      \\[0.2cm]
      we recursively evaluate the expressions \texttt{lhs} and \texttt{rhs} and return the sum of these
      values. 
\item The evaluation of the arithmetic operators \texttt{'-'}, \texttt{'*'}, and \texttt{'/'}
      is analogous to the evaluation of the operator \texttt{'+'}.
\end{enumerate}
\FloatBarrier

\noindent
The Jupyter notebook \texttt{03-Interpreter.ipynb}, which is available at
\\[0.2cm]
\hspace*{0.3cm}
\href{https://github.com/karlstroetmann/Formal-Languages/blob/master/Python/Chapter-07/03-Interpreter.ipynb}{https://github.com/karlstroetmann/Formal-Languages/blob/master/Python/Chapter-07/03-Interpreter.ipynb}
\\[0.2cm]
implements an interpreter.

\exerciseEng
\begin{enumerate}[(a)]
\item Add \texttt{for} loops to the interpreter.
\item Expand the interpreter to include logical operators
      ``\texttt{\&\&}'' for logical \emph{and}, ``\texttt{||}'' for logical \emph{or},
      and ``\texttt{!}'' for \emph{negation}. The operator ``\texttt{!}'' should bind 
      strongest and the operator ``\texttt{||}'' should bind weakest.
\item Enhance the interpreter to allow for user-defined functions.
      \begin{itemize}
      \item A function should only have access to its parameters and those variables that
            are defined locally inside the function.  
      \item A function should always return a value using a \texttt{return} statement.
            In order to facilitate that a function can contain multiple \texttt{return}
            statements and that a \texttt{return} statement can occur
            anywhere in the function,  use an \blue{exception} to communicate the
            return value and to transfer the control flow out of the function.
      \end{itemize}
\end{enumerate}
Figure \ref{fig:factorial.sl} on page \pageref{fig:factorial.sl} shows a simple program
that implements a function to compute the factorial function.  Your interpreter should be able to execute this
program.
Use the Jupyter notebook \texttt{Interpreter-Frame.ipynb}, which is available at
\\[0.2cm]
\hspace*{0.3cm}
\href{https://github.com/karlstroetmann/Formal-Languages/blob/master/Python/Chapter-07/Interpreter-Frame.ipynb}{https://github.com/karlstroetmann/Formal-Languages/blob/master/Python/Chapter-07/Interpreter-Frame.ipynb}
\\[0.2cm]
as a starting point for your solution.
\eox

\begin{figure}[!ht]
\centering
\begin{minted}[ frame         = lines, 
                 framesep      = 0.3cm, 
                 firstnumber   = 1,
                 bgcolor       = sepia,
                 numbers       = left,
                 numbersep     = -0.2cm,
                 xleftmargin   = 0.8cm,
                 xrightmargin  = 0.8cm,
                 ]{javascript}
    function factorial(n) {
        if (n == 0) {
            return 1;
        }
        return n * factorial(n - 1);
    }
    
    for (i := 2; i <= 25; i := i + 1) {
        print(factorial(i));
    }                 
\end{minted}
\vspace*{-0.3cm}
\caption{A program to compute the factorial function.}
\label{fig:factorial.sl}
\end{figure}
\pagebreak

\section{Check your Understanding}
\begin{enumerate}[(a)]
\item Are you able to implement a simple calculator using \simtextsc{PLy}?
\item Would you be able to implement a \simtextsc{Ply} parser that reads propositional formulas and returns these
      formulas as nested tuples? 
\item Are you able to extend the interpreter that has been discussed in this chapter?
\item Assume that we develop an interpreter that supports user defined functions.
      How can we accomodate the fact that a \texttt{return} statement can occur anywhere in a function? 
\end{enumerate}



%%% Local Variables: 
%%% mode: latex
%%% TeX-master: "formal-languages"
%%% End: 
